\section{Conclusion}
\label{sec:conclusion}

Pakistan's commercial aviation fuel and OECD flight estimates nearly close in aggregate in the
latest two fiscal years, but that agreement is produced by large component differences of
opposite sign. In FY2023--24 resident-customer fuel is 28.6\% above the resident-carrier
estimate, foreign-customer fuel is 23.1\% below the non-resident estimate and 97\% of the
discrepancy cancels. Pakistan's official 80/20 allocation produces a still larger domestic--
international mismatch. National agreement is therefore a necessary gross check, not evidence
that purchaser categories implement flight geography.

The national inventory carries the aggregate movement model's \SI{0.301}{\mega\tonne}
\COe{} as domestic aviation's figure. The central bridge assigns 0.292 Mt \COtwo{};
aggregate-movement, route-resolved and completed route--aircraft estimates are 0.299,
0.271 and 0.247 Mt as \COtwo{}; the route--aircraft Monte Carlo mean is 0.244 Mt with a
95\% interval of 0.215--0.272 Mt. These estimates place domestic aviation in a common broad range while showing
that customer partition, aircraft assignment and incomplete international coverage remain the
most consequential uncertainties.

The general contribution is an accounting method rather than one assertedly final national
number. Every fuel source row is conserved, residuals remain explicit, evidence roles are
declared and structural mappings are kept outside parametric confidence intervals. Observed
flight stages and detailed uplift records can now test and replace
provisional inputs without changing those rules.
